\section{固体の内部で局所的に融解する}\label{theory:volcano:melting}
マントルの多くは固体である。地震波を伝える固体であることと、地質学的な時間にゆっくり流動することは両立する。マグマは地球内部を一面に覆う液体ではなく、温度・圧力・組成の経路が融解条件を越える場所で生じる。
多成分岩石では融解開始の固相線と全てが融ける液相線が異なる。間では固相とメルトが共存し、部分溶融となる。初生マグマは原岩の組成をそのまま全量コピーしたものではない\cite{vol50_melt_ref}。
物理での相図の読み方を使うと、融解開始は概念的に \(T_{\rm path}(P)>T_{\rm solidus}(P,X)\) と表せる。上昇するマントルでは圧力が下がり、ほぼ断熱的な温度経路と固相線の交差で減圧融解が始まる。海嶺の成因を「低温の物質が突然加熱される」と説明しない。ホットスポットの熱異常と上昇減圧は別の要因である。
沈み込みでは含水物質の脱水により流体がマントルウェッジへ供給され、融解開始温度を下げる\cite{vol50_subduction_ref}。温度を上げる、圧力を下げる、水等で固相線を下げるという三経路を分ける。
メルトが周囲より低密度なら浮力で上昇し得る。概略の駆動力は体積 \(V\) 当たり \((\rho_{\rm rock}-\rho_{\rm melt})gV\) だが、実際の移動には割れ目・応力・粘性・冷却も関係する。全マグマだまりを中立浮力だけで決定しない。

\section{原岩と別マグマのどちらと比較するか}
初生的なマントル融解液は代表的に玄武岩質である。結晶が分離すれば残る液体の組成は変わり、地殻の融解や同化・混合も別のマグマを作る\cite{vol50_melt_ref}。
玄武岩質は流紋岩質等よりSiO2に乏しい側だが、原岩マントルに比べればNa・K等の不適合元素がメルトへ入り濃集する場合がある。「富む」は比較対象を明記する。
JMAの1999年表では玄武岩質の代表噴出温度が1000--1200℃、流紋岩質が700--900℃で、珪酸分が多い側は一般に粘性が大きい\cite{jma_magma1999}。これは原典時点の分類と代表値であり、全組成・全含水量・全結晶量の固定境界ではない。

\section{海溝・前弧・火山弧・背弧}\label{theory:volcano:arc}
地表で海溝から陸側へ前弧・火山弧・背弧を並べ、地下では沈み込むスラブとその上のマントルウェッジを対応させる。火山前線は火山の分布の海溝側限界で、海溝そのものと違う\cite{usgs_plate_boundaries}。
火山が少ない前弧でもプレート境界やスラブの地震活動があり、「火山がないから地震もない」は誤りである。背弧に拡大海盆ができる場合はあるが、全島弧が同じ拡大状態を持つとはしない。
アルカリ元素と水を一つの変数として扱わない。島弧の背弧方向のアルカリ濃度には一般的な増加傾向を学ぶ一方、GSJの東北日本弧の発表には含水量が背弧へ減少する分布の報告がある\cite{vol50_arcwater_ref}。融解度・供給源・脱ガスの違いを考え、全弧の水濃度へ固定した順位を使わない。
中央海嶺は発散境界で新しい海洋地殻を作る。USGS概説では毎年の地球の噴出物の約4分の3が拡大海嶺起源と推計される\cite{vol50_ridge_ref}。陸上噴火の個数や被害の印象と、全世界の体積率を区別する。

\section{間欠爆発と連続溶岩噴泉}\label{theory:volcano:styles}
低粘性の玄武岩質マグマの連続的な噴泉や流出はハワイ式の代表である。ストロンボリ式はガスの断続的な解放により小爆発を繰り返し、火山弾・スコリアを放出する\cite{jma_magma1999,vol50_styles_ref}。同じ火山が常に一つの様式だけで噴火するわけではない。
高粘性の溶岩は流れにくく火口近くへ蓄積してドームを作り得る。ガスが抜ける時間、供給量、結晶量等も噴火様式へ関係する。岩種名だけで爆発の有無を必ず決めない。
マグマ水蒸気爆発では外来水とマグマの急な熱交換・蒸気膨張・破砕が関与する。水蒸気噴火で既存岩石が破砕される場合と、新しいマグマが外来水へ接触する場合を区別する。

\section{溶岩表面の形と冷却}\label{theory:volcano:lava}
パホイホイは代表的に玄武岩質の流動的な溶岩で、滑らか・縄状等の表面を持つ。アアは粗いクリンカーに覆われる。同じ玄武岩でも冷却・結晶化・変形速度等で一方から他方へ変化し得る\cite{vol50_lava_ref}。滑らかだから流紋岩という対応は誤りである。
ブロック溶岩では大きい岩塊の比較的平滑な破断面が特徴となる。表面を形態で分ける軸と、岩石の化学組成を分ける軸を混同しない。
厚さ／長さ等のアスペクト比は地形を比較する指標である。原典の1/8区分は分類上の目安として保存し、普遍的な相転移や流動停止の定数とはしない。長さの測り方・流量・地形・冷却の影響を考える。

\section{ガスの主要成分と平衡条件}\label{theory:volcano:gases}
通常の火山ガスの最多成分は水蒸気であり、CO2、SO2、H2S、水素ハロゲン化物等も放出される\cite{vol50_gas_ref}。危険性の大きさと濃度順位は別の情報である。
減圧ではメルト中の揮発性分の溶解度が下がり、気泡が分離し膨張する。CO2と水の脱ガス段階は同じではない\cite{vol50_degassing_ref}。閉塞で過圧が生じる場合と、開いた系でガスが抜ける場合を分ける。外圧低下だけからガス絶対圧力の単調増加を推論しない。
高温マグマのガス、冷却された噴気、地下水と反応した熱水ガスは組成が違い得る。CH4を通常高温の主要ガスの比較で選外とすることは、全自然条件で存在しないことを意味しない。令和7年度の100気圧・1200℃・玄武岩質という条件も、酸化還元状態・全元素量なしに任意の正確なモル割合を作れるわけではない。

\section{粒子とガス、水、岩屑を分ける}\label{theory:volcano:hazards}
火砕流は高温のガス・粒子混合、火山泥流は水・粒子混合、岩屑なだれは崩壊岩体の移動を中心に考える。速さの目安と材料を同時に確認する。噴煙柱の崩壊やドーム崩落と、融雪・湖水・降雨による泥流形成を同一視しない。火砕密度流の背景は第\ref{theory:volcanology:pyroflow}節へ戻る。
大量の物質が海へ入り水を急激に動かすと津波が生じる場合がある。噴火・海底地殻変動・崩壊による水塊の変位を考え、全噴火が津波を起こすとはしない。
噴煙が高くても、灰の短期的な遮蔽とSO2由来硫酸エアロゾルの長距離・長期間の気候影響を区別する\cite{vol50_climate_ref}。粒子の沈降時間、化学反応、成層圏での滞留という物理と化学の段階を追う。

\section{火炎状構造の「火」は成因ではない}\label{theory:volcano:terms}
カルデラ、火山岩の安山岩、溶岩片が付着したアグルチネート、マグマ水蒸気爆発は火山の直接用語である。
火炎状構造は未固結の軽い堆積層が上の重い層の間へ押し上げられる等の変形を指し、名称は形の比喩である\cite{vol50_terms_ref}。堆積物の形成地域が火山である可能性を認めても、変形の直接機構を炎や噴火としない。名前を覚えた後に、物質・応力・流れの機構を説明できることを学習目標とする。
